% =============================================================================
% APPENDIX AZ: METHOD-CONSTRUCT: Methodological Foundations of Model Construction
% =============================================================================
% Module: BCM2_08_META (Methodology)
% Version: 2.0 (January 2026)
% Status: Review
%
% This appendix addresses the methodological foundations required for rigorous
% model construction in EBF: microfoundations, learning theory, dynamics,
% heterogeneity, and model identification.
%
% =============================================================================

% =============================================================================
% CROSS-REFERENCE STRUCTURE
% =============================================================================
%
% CHAPTER LINKAGE (Required):
% - Primary Chapter: Chapter 8: Mathematical Foundations
% - Secondary Chapters: Chapter 4x (Calibration), Chapter 2 (Rationality)
%
% DEPENDENCIES (what this appendix REQUIRES):
% - Appendix A (FORMAL-DERIVE): Mathematical foundations
% - Appendix D (FORMAL-PROOF): C* derivations and proofs
% - Appendix T (REF-META): Metatheoretical foundations
% - Appendix BBB (CORE-WHERE): Parameter estimation framework
%
% DEPENDENTS (what REQUIRES this appendix):
% - All CORE appendices: Methodological grounding
% - Appendix AN (METHOD-LLMMC): Estimation methodology
% - Appendix R (METHOD-EVAL): Validation framework
%
% =============================================================================

\begin{tcolorbox}[colback=purple!5!white, colframe=purple!75!black,
    title=Chapter Linkage]

\textbf{Primary Chapter:} Chapter 8: Mathematical Foundations
\begin{itemize}[nosep]
\item Section 8.1: Formal Framework $\leftrightarrow$ This Appendix Section \ref{sec:az-micro}
\item Section 8.2: Equilibrium Concepts $\leftrightarrow$ This Appendix Section \ref{sec:az-dynamics}
\end{itemize}

\textbf{Secondary Chapters:}
\begin{itemize}[nosep]
\item Chapter 4x: Calibration not Simulation --- LLM MC estimation methodology
\item Chapter 2: Rationality as Stability --- stability-based reference structure
\item Chapter 6: Reference Structure --- $C^*$ fixed-point theory
\end{itemize}

\textbf{Reading Guide:}
\begin{itemize}[nosep]
\item For conceptual overview: Read Chapter 8 first
\item For formal C* theory: See Appendix D (FORMAL-PROOF)
\item For estimation methods: See Appendix AN (METHOD-LLMMC)
\end{itemize}

\end{tcolorbox}

\chapter{METHOD-CONSTRUCT: Methodological Foundations of Model Construction}
\label{app:az-construct}

\begin{abstract}
This appendix addresses five fundamental methodological challenges in constructing
EBF: (1) Microfoundations---how axioms derive from primitive assumptions;
(2) Learning Theory---how agents form and update expectations;
(3) Dynamic Adjustment---how systems converge to reference structure $C^*$;
(4) Heterogeneity and Aggregation---how individual differences scale;
(5) Identification and Model Comparison---how to distinguish EBF from alternatives.
These foundations ensure EBF meets rigorous standards of economic methodology.
\end{abstract}

\begin{tcolorbox}[colback=blue!5!white, colframe=blue!75!black,
    title=Quick Reference: METHOD-CONSTRUCT]

\textbf{Core Question:} What methodological foundations make EBF a rigorous scientific model?

\textbf{Key Formula:}
\begin{equation}
C^* = \Phi(C^*, \Psi) \quad \text{where} \quad \Phi: \mathcal{C} \times \mathcal{P} \to \mathcal{C}
\end{equation}

\textbf{Key Concepts:}
\begin{itemize}[nosep]
\item \textbf{Microfoundations}: Deriving macro behavior from individual axioms
\item \textbf{Learning Dynamics}: Bayesian updating toward $C^*$
\item \textbf{Identification}: Distinguishing EBF from nested alternatives
\item \textbf{Aggregation}: From heterogeneous agents to population behavior
\item \textbf{Convergence}: Conditions for dynamic stability of $C^*$
\end{itemize}

\textbf{Cross-References:} Appendix A (FORMAL-DERIVE), Appendix D (FORMAL-PROOF),
Appendix T (REF-META), Appendix BBB (CORE-WHERE), Appendix AN (METHOD-LLMMC)
\end{tcolorbox}


% =============================================================================
%                    PART 2: CORE CONTENT
% =============================================================================

%%%%%%%%%%%%%%%%%%%%%%%%%%%%%%%%%%%%%%%%%%%%%%%%%%%%%%%%%%%%%%%%%%%%%%%%%%%%%%%
\section{The Fundamental Question}
\label{sec:az-fundamental}
%%%%%%%%%%%%%%%%%%%%%%%%%%%%%%%%%%%%%%%%%%%%%%%%%%%%%%%%%%%%%%%%%%%%%%%%%%%%%%%

\subsection{Why This Matters}

Economic models face five persistent methodological challenges that determine
their scientific credibility:

\begin{enumerate}
\item \textbf{Microfoundation Critique}: Models must derive aggregate behavior
from individual decision-making, not impose ad hoc functional forms
\citep{lucas1976econometric}.

\item \textbf{Learning Problem}: Rational expectations require agents to know
model parameters---but how do they learn them? \citep{sargent1993bounded}.

\item \textbf{Dynamics Gap}: Static equilibria ignore adjustment processes and
out-of-equilibrium behavior \citep{arthur2021foundations}.

\item \textbf{Aggregation Problem}: How do heterogeneous individual behaviors
combine into population-level patterns? \citep{kirman1992whom}.

\item \textbf{Identification Crisis}: Many models generate identical predictions,
making empirical discrimination impossible \citep{manski1995identification}.
\end{enumerate}

EBF must address each challenge explicitly to claim scientific rigor.

\subsection{The Core Problem}

\begin{tcolorbox}[colback=red!5!white, colframe=red!75!black,
    title=The Fundamental Question]
\textbf{How does EBF establish methodological foundations that satisfy
the standards of microfoundation, learning consistency, dynamic stability,
valid aggregation, and empirical identification?}
\end{tcolorbox}


%%%%%%%%%%%%%%%%%%%%%%%%%%%%%%%%%%%%%%%%%%%%%%%%%%%%%%%%%%%%%%%%%%%%%%%%%%%%%%%
\section{Part 1: Microfoundations of the Axiom System}
\label{sec:az-micro}
%%%%%%%%%%%%%%%%%%%%%%%%%%%%%%%%%%%%%%%%%%%%%%%%%%%%%%%%%%%%%%%%%%%%%%%%%%%%%%%

\subsection{The Microfoundation Requirement}

Following \citet{lucas1976econometric}, aggregate behavior must emerge from
individual optimization. EBF satisfies this through:

\begin{tcolorbox}[colback=gray!5!white, colframe=gray!75!black,
    title=Axiom AZ-M1: Methodological Individualism with Levels]
\textbf{Statement:} All aggregate welfare $U^L$ for $L > 0$ is reducible to
weighted combinations of individual welfare $U^0_i$:
\begin{equation}
U^L = \sum_{i \in M^L} \omega^L_i \cdot U^0_i + \sum_{i < j} \gamma^L_{ij} \cdot U^0_i \cdot U^0_j
\end{equation}

\textbf{Interpretation:} Higher-level utility exists only through individuals who
experience it. The complementarity term $\gamma^L_{ij}$ captures genuine emergence.

\textbf{Implication:} No ``social organism'' fallacy---levels are epistemological,
not ontological.
\end{tcolorbox}

\subsection{Deriving Axioms from Primitives}

The EBF axiom system derives from four primitive assumptions:

\begin{enumerate}
\item \textbf{P1 (Preference Existence)}: Agent $i$ has preferences $\succeq_i$
over outcome bundles $\mathbf{x} \in \mathcal{X}$.

\item \textbf{P2 (Completeness)}: For all $\mathbf{x}, \mathbf{y} \in \mathcal{X}$:
$\mathbf{x} \succeq_i \mathbf{y}$ or $\mathbf{y} \succeq_i \mathbf{x}$.

\item \textbf{P3 (Transitivity)}: For all $\mathbf{x}, \mathbf{y}, \mathbf{z} \in \mathcal{X}$:
if $\mathbf{x} \succeq_i \mathbf{y}$ and $\mathbf{y} \succeq_i \mathbf{z}$, then
$\mathbf{x} \succeq_i \mathbf{z}$.

\item \textbf{P4 (Context Sensitivity)}: Preferences are context-indexed:
$\succeq_i(\Psi)$ may differ from $\succeq_i(\Psi')$ for $\Psi \neq \Psi'$.
\end{enumerate}

\begin{proposition}[Representation Theorem]
Under P1--P4, there exists a utility function $U^i: \mathcal{X} \times \mathcal{P} \to \mathbb{R}$
such that:
\begin{equation}
\mathbf{x} \succeq_i(\Psi) \mathbf{y} \Leftrightarrow U^i(\mathbf{x}, \Psi) \geq U^i(\mathbf{y}, \Psi)
\end{equation}
\end{proposition}

\begin{proof}
By Debreu's Representation Theorem applied separately for each $\Psi$, with
measurable dependence on $\Psi$ ensured by P4's regularity conditions.
See Appendix D, Section D.1 for complete proof.
\end{proof}

\subsection{From Primitives to FEPSDE}

The six FEPSDE dimensions arise from decomposing the outcome space:

\begin{equation}
\mathcal{X} = \mathcal{X}_F \times \mathcal{X}_E \times \mathcal{X}_P \times
\mathcal{X}_S \times \mathcal{X}_D \times \mathcal{X}_{Ec}
\end{equation}

\begin{tcolorbox}[colback=gray!5!white, colframe=gray!75!black,
    title=Axiom AZ-M2: Dimensional Separability]
\textbf{Statement:} Utility admits additive separable representation across dimensions:
\begin{equation}
U^i(\mathbf{x}, \Psi) = \sum_{d \in FEPSDE} U^i_d(x_d, \Psi) + \sum_{d < d'} \gamma_{dd'}(\Psi) \cdot U^i_d \cdot U^i_{d'}
\end{equation}

\textbf{Interpretation:} Dimensions are conceptually distinct but may interact
through complementarities $\gamma_{dd'}$.

\textbf{Testable Prediction:} Cross-partial derivatives $\partial^2 U / \partial x_d \partial x_{d'}$
should have consistent sign structure.
\end{tcolorbox}


%%%%%%%%%%%%%%%%%%%%%%%%%%%%%%%%%%%%%%%%%%%%%%%%%%%%%%%%%%%%%%%%%%%%%%%%%%%%%%%
\section{Part 2: Learning Theory and Expectation Formation}
\label{sec:az-learning}
%%%%%%%%%%%%%%%%%%%%%%%%%%%%%%%%%%%%%%%%%%%%%%%%%%%%%%%%%%%%%%%%%%%%%%%%%%%%%%%

\subsection{The Learning Problem}

Rational expectations models assume agents know the true model parameters
$\Theta$. But in reality, agents must \textit{learn} these parameters from
experience \citep{sargent1993bounded, evans2001learning}.

EBF addresses this through a Bayesian learning framework:

\begin{tcolorbox}[colback=gray!5!white, colframe=gray!75!black,
    title=Axiom AZ-L1: Bayesian Parameter Learning]
\textbf{Statement:} Agent $i$ updates beliefs about parameters $\Theta$ according to:
\begin{equation}
p(\Theta | \mathcal{H}_t) \propto p(\mathcal{H}_t | \Theta) \cdot p(\Theta | \mathcal{H}_{t-1})
\end{equation}
where $\mathcal{H}_t$ denotes the agent's information history up to time $t$.

\textbf{Interpretation:} Agents are Bayesian learners who update parameter beliefs
based on observed outcomes.

\textbf{Implication:} Prior beliefs matter for behavior, especially in novel contexts.
\end{tcolorbox}

\subsection{Learning the Reference Structure $C^*$}

The endogenous reference structure $C^*$ must itself be learned. This creates
a fixed-point learning problem:

\begin{equation}
C^*_t = \mathbb{E}[\Phi(C^*_t, \Psi) | \mathcal{H}_t]
\end{equation}

\begin{proposition}[Adaptive Learning Convergence]
Under regularity conditions (Lipschitz $\Phi$, bounded noise, persistent excitation),
the adaptive learning process:
\begin{equation}
C_{t+1} = C_t + \alpha_t \cdot (\Phi(C_t, \Psi_t) - C_t) + \epsilon_t
\end{equation}
converges almost surely to $C^*$ as $t \to \infty$, where $\sum_t \alpha_t = \infty$
and $\sum_t \alpha_t^2 < \infty$.
\end{proposition}

\begin{proof}
By stochastic approximation theory \citep{benveniste1990adaptive}.
See Appendix D, Section D.2 for complete conditions.
\end{proof}

\subsection{Expectation Formation under Bounded Rationality}

EBF allows for bounded rational expectation formation:

\begin{tcolorbox}[colback=gray!5!white, colframe=gray!75!black,
    title=Axiom AZ-L2: Hierarchical Expectations]
\textbf{Statement:} Agent $i$'s expectations are formed at cognitive level $k$:
\begin{equation}
\mathbb{E}^{(k)}_i[X] = \begin{cases}
\bar{X} & k = 0 \text{ (naive)} \\
\mathbb{E}[X | \mathcal{H}_i] & k = 1 \text{ (adaptive)} \\
\mathbb{E}[X | \Theta^*, \mathcal{H}_i] & k = 2 \text{ (rational)} \\
\mathbb{E}[\mathbb{E}^{(k-1)}_j[X]] & k \geq 3 \text{ (higher-order)}
\end{cases}
\end{equation}

\textbf{Interpretation:} Different cognitive levels yield different expectation
formation processes. The Awareness function $A(\cdot)$ determines which level is used.

\textbf{Implication:} Connects to level-$k$ thinking literature \citep{stahl1995players,
nagel1995unraveling}.
\end{tcolorbox}


%%%%%%%%%%%%%%%%%%%%%%%%%%%%%%%%%%%%%%%%%%%%%%%%%%%%%%%%%%%%%%%%%%%%%%%%%%%%%%%
\section{Part 3: Dynamic Adjustment to $C^*$}
\label{sec:az-dynamics}
%%%%%%%%%%%%%%%%%%%%%%%%%%%%%%%%%%%%%%%%%%%%%%%%%%%%%%%%%%%%%%%%%%%%%%%%%%%%%%%

\subsection{The Dynamics Problem}

Static equilibrium analysis ignores the process by which systems reach equilibrium.
EBF specifies explicit adjustment dynamics:

\begin{tcolorbox}[colback=gray!5!white, colframe=gray!75!black,
    title=Axiom AZ-D1: Gradient Adjustment Dynamics]
\textbf{Statement:} The reference structure evolves according to:
\begin{equation}
\frac{dC}{dt} = \eta \cdot \nabla_C V(C, \Psi) = \eta \cdot (\Phi(C, \Psi) - C)
\end{equation}
where $V(C, \Psi)$ is a potential function and $\eta > 0$ is adjustment speed.

\textbf{Interpretation:} The system moves in the direction of increasing coherence $K$.

\textbf{Implication:} $C^*$ is a stable attractor when $\nabla^2_C V(C^*, \Psi) < 0$.
\end{tcolorbox}

\subsection{Stability Analysis}

The stability of $C^*$ depends on the eigenvalues of the Jacobian:

\begin{equation}
J = \frac{\partial \Phi}{\partial C}\bigg|_{C=C^*}
\end{equation}

\begin{proposition}[Local Stability]
$C^*$ is locally asymptotically stable if and only if all eigenvalues of $J$
have modulus less than 1 (discrete time) or negative real part (continuous time).
\end{proposition}

\begin{proof}
Standard linearization argument. See Appendix D, Section D.3.
\end{proof}

\subsection{Multiple Equilibria and Path Dependence}

The mapping $\Phi$ may have multiple fixed points:

\begin{tcolorbox}[colback=yellow!5!white, colframe=yellow!75!black,
    title=Result AZ-D2: Multiplicity of $C^*$]
When $\Phi$ is non-monotonic, multiple stable $C^*$ may exist:
\begin{equation}
\{C^*_1, C^*_2, \ldots, C^*_n\} = \{C : \Phi(C, \Psi) = C, \text{ locally stable}\}
\end{equation}

\textbf{Implications:}
\begin{itemize}[nosep]
\item History matters: Initial conditions determine which $C^*$ is reached
\item Hysteresis: Small shocks may not move system, large shocks may cause regime shift
\item Policy sensitivity: Interventions may have permanent effects on reference structure
\end{itemize}
\end{tcolorbox}

\subsection{Speed of Convergence}

The rate of convergence to $C^*$ is:

\begin{equation}
\|C_t - C^*\| \approx \|C_0 - C^*\| \cdot e^{-\lambda_{\min} t}
\end{equation}

where $\lambda_{\min}$ is the smallest eigenvalue magnitude of $I - J$.

\textbf{Empirical Prediction:} Reference points adjust slowly (half-life of months
to years) rather than instantaneously.


%%%%%%%%%%%%%%%%%%%%%%%%%%%%%%%%%%%%%%%%%%%%%%%%%%%%%%%%%%%%%%%%%%%%%%%%%%%%%%%
\section{Part 4: Heterogeneity and Aggregation}
\label{sec:az-heterogeneity}
%%%%%%%%%%%%%%%%%%%%%%%%%%%%%%%%%%%%%%%%%%%%%%%%%%%%%%%%%%%%%%%%%%%%%%%%%%%%%%%

\subsection{The Aggregation Problem}

Individual heterogeneity poses fundamental challenges for aggregation
\citep{kirman1992whom}. EBF addresses this through:

\begin{tcolorbox}[colback=gray!5!white, colframe=gray!75!black,
    title=Axiom AZ-H1: Distributional Heterogeneity]
\textbf{Statement:} Individual parameters follow distributions:
\begin{equation}
\Theta_i \sim F(\Theta | \Psi) \quad \text{where} \quad \Theta_i = \{\gamma_i, \omega_i, \lambda_i, \delta_i, A_i, \theta_i\}
\end{equation}

\textbf{Interpretation:} Agents differ in complementarity, weights, loss aversion,
time preferences, awareness, and action thresholds.

\textbf{Implication:} Population behavior reflects the entire distribution, not just means.
\end{tcolorbox}

\subsection{Aggregation Rules}

Aggregate behavior emerges through:

\begin{equation}
\bar{U}^L(\Psi) = \int U^L_i(\Psi; \Theta_i) \, dF(\Theta_i | \Psi)
\end{equation}

\begin{proposition}[Representative Agent Conditions]
A representative agent with parameters $\bar{\Theta}$ exists if and only if:
\begin{enumerate}
\item Utility is quasi-linear in at least one dimension, or
\item Heterogeneity is restricted to scale (homothetic preferences), or
\item Aggregation is over a single level $L$
\end{enumerate}
\end{proposition}

\begin{proof}
By Gorman aggregation theorem \citep{gorman1953community}.
See Appendix D, Section D.4.
\end{proof}

\subsection{Beyond Representative Agent}

When representative agent conditions fail, EBF uses distributional aggregation:

\begin{equation}
P(\text{Action} | \Psi) = \int \mathbf{1}[WAX_i(\Psi; \Theta_i) \geq \theta_i] \, dF(\Theta_i | \Psi)
\end{equation}

This yields predictions about population shares rather than aggregate behavior.

\subsection{Heterogeneity Types in EBF}

\begin{table}[htbp]
\centering
\caption{Sources of Heterogeneity in EBF}
\label{tab:az-heterogeneity}
\renewcommand{\arraystretch}{1.3}
\begin{tabular}{@{}llll@{}}
\toprule
\textbf{Parameter} & \textbf{Symbol} & \textbf{Source} & \textbf{Effect} \\
\midrule
Complementarity & $\gamma_i$ & Cognitive style & Synergy perception \\
Dimension weights & $\omega_{id}$ & Values, culture & Priority structure \\
Loss aversion & $\lambda_i$ & Experience, personality & Risk response \\
Time preference & $\delta_i$ & Age, wealth & Intertemporal tradeoffs \\
Awareness & $A_i(\cdot)$ & Attention, salience & What matters at $t^*$ \\
Threshold & $\theta_i$ & Habits, friction & Action barrier \\
Thinking style & $\varphi_i$ & Cognitive mode & Additive vs. complementary \\
\bottomrule
\end{tabular}
\end{table}

\subsection{Empirical Parameter Estimates from Literature}

The following table summarizes empirical estimates for key EBF parameters
from the behavioral economics literature:

\begin{table}[htbp]
\centering
\caption{Empirical Estimates of EBF Parameters}
\label{tab:az-empirical}
\renewcommand{\arraystretch}{1.3}
\small
\begin{tabular}{@{}llllp{4cm}@{}}
\toprule
\textbf{Parameter} & \textbf{Estimate} & \textbf{95\% CI} & \textbf{Source} & \textbf{Context} \\
\midrule
\multicolumn{5}{@{}l}{\textit{Learning and Adjustment}} \\
$\alpha$ (learning rate) & 0.05--0.15 & [0.02, 0.25] & Evans \& Honkapohja & Macro expectations \\
$\eta$ (adjustment speed) & 0.08 & [0.05, 0.12] & This appendix & Carbon tax response \\
$t_{1/2}$ (half-life) & 6--12 mo & [3, 24] & Various & Reference adjustment \\
\addlinespace
\multicolumn{5}{@{}l}{\textit{Loss Aversion and Risk}} \\
$\lambda$ (loss aversion) & 2.25 & [1.5, 3.0] & Tversky \& Kahneman & Prospect theory \\
$\lambda_F$ (financial) & 2.0 & [1.5, 2.5] & Abdellaoui et al. & Financial decisions \\
$\lambda_S$ (social) & 2.8 & [2.0, 3.5] & Baumeister et al. & Social rejection \\
\addlinespace
\multicolumn{5}{@{}l}{\textit{Time Preferences}} \\
$\delta$ (discount factor) & 0.90 & [0.80, 0.98] & Frederick et al. & Annual rate \\
$\beta$ (present bias) & 0.70 & [0.5, 0.9] & Laibson & Quasi-hyperbolic \\
\addlinespace
\multicolumn{5}{@{}l}{\textit{Social Preferences}} \\
$\gamma_{ij}$ (altruism) & 0.30 & [0.1, 0.5] & Fehr \& Schmidt & Dictator games \\
$\omega_{family}$ & 0.65 & [0.4, 0.8] & Becker & Family utility weight \\
\addlinespace
\multicolumn{5}{@{}l}{\textit{Awareness and Thresholds}} \\
$A$ (mean awareness) & 0.40 & [0.2, 0.6] & Bordalo et al. & Salience theory \\
$\theta$ (action threshold) & 0.55 & [0.3, 0.8] & Various & Intention-behavior gap \\
$\varphi$ (thinking style) & 0.45 & [0.2, 0.7] & Toplak et al. & CRT distribution \\
\bottomrule
\end{tabular}
\end{table}

\textbf{Interpretation Notes:}
\begin{itemize}[nosep]
\item These estimates are \textit{illustrative starting points}---actual values
vary substantially by context $\Psi$.
\item Epistemic status: Most are THR (theoretical) or EMP (empirical); see
Appendix BBB for the complete hierarchy.
\item Heterogeneity: Standard deviations are typically 30--50\% of the mean,
indicating substantial individual differences.
\end{itemize}


%%%%%%%%%%%%%%%%%%%%%%%%%%%%%%%%%%%%%%%%%%%%%%%%%%%%%%%%%%%%%%%%%%%%%%%%%%%%%%%
\section{Part 5: Identification and Model Comparison}
\label{sec:az-identification}
%%%%%%%%%%%%%%%%%%%%%%%%%%%%%%%%%%%%%%%%%%%%%%%%%%%%%%%%%%%%%%%%%%%%%%%%%%%%%%%

\subsection{The Identification Challenge}

Model identification requires that EBF makes predictions distinguishable
from alternatives \citep{manski1995identification}:

\begin{tcolorbox}[colback=gray!5!white, colframe=gray!75!black,
    title=Axiom AZ-I1: Falsifiability Requirement]
\textbf{Statement:} For every core parameter $\theta \in \Theta$, there exists
an observable implication $\phi(\theta)$ such that:
\begin{equation}
\theta \neq \theta' \Rightarrow \phi(\theta) \neq \phi(\theta')
\end{equation}

\textbf{Interpretation:} Different parameter values must yield different observable predictions.

\textbf{Implication:} EBF parameters are empirically meaningful, not arbitrary.
\end{tcolorbox}

\subsection{Nested Model Comparisons}

EBF nests several standard models as special cases:

\begin{table}[htbp]
\centering
\caption{EBF Nesting Structure}
\label{tab:az-nesting}
\renewcommand{\arraystretch}{1.3}
\begin{tabular}{@{}lll@{}}
\toprule
\textbf{Nested Model} & \textbf{Restriction} & \textbf{Test Statistic} \\
\midrule
Standard EU & $\gamma = 0$, $A = 1$, $L = 0$ & LR test on $\gamma$ \\
Prospect Theory & $\gamma = 0$, $A = 1$, $\lambda \neq 1$ & LR test on $\gamma$, $L$ \\
Social Preferences & $L > 0$, $\gamma = 0$ & LR test on $\gamma$ \\
Intertemporal Choice & Multi-$t$, $\gamma = 0$ & Cross-dimension $\gamma$ \\
Behavioral: all levels & All $L$, $\gamma \neq 0$, $A < 1$ & Full EBF \\
\bottomrule
\end{tabular}
\end{table}

\subsection{Identification Strategies}

\begin{enumerate}
\item \textbf{Within-Person Variation:} Same person, different contexts $\Psi$
\begin{equation}
\Delta U_i = U_i(\Psi') - U_i(\Psi) \Rightarrow \text{identifies } \Gamma_i(\Psi)
\end{equation}

\item \textbf{Cross-Level Variation:} Same context, different aggregation levels
\begin{equation}
U^{L'} - U^L \Rightarrow \text{identifies } \gamma^{cross}_{L,L'}
\end{equation}

\item \textbf{Temporal Variation:} Same person and context, different time points
\begin{equation}
C^*_{t'} - C^*_t \Rightarrow \text{identifies adjustment dynamics}
\end{equation}

\item \textbf{Experimental Variation:} Randomized interventions
\begin{equation}
U(\text{treatment}) - U(\text{control}) \Rightarrow \text{identifies causal effects}
\end{equation}
\end{enumerate}

\subsection{Model Selection Criteria}

\begin{tcolorbox}[colback=yellow!5!white, colframe=yellow!75!black,
    title=Result AZ-I2: Model Selection Protocol]
Compare EBF to alternatives using:
\begin{enumerate}
\item \textbf{In-sample fit:} $R^2$, likelihood, RMSE
\item \textbf{Out-of-sample prediction:} Cross-validation, holdout sets
\item \textbf{Parsimony:} AIC, BIC penalizing parameters
\item \textbf{Theoretical coherence:} Axiom consistency
\item \textbf{Novel predictions:} Predictions unique to EBF
\end{enumerate}

EBF is preferred when it provides superior out-of-sample prediction
while satisfying axiom consistency and generating verified novel predictions.
\end{tcolorbox}


%%%%%%%%%%%%%%%%%%%%%%%%%%%%%%%%%%%%%%%%%%%%%%%%%%%%%%%%%%%%%%%%%%%%%%%%%%%%%%%
\section{Part 6: Computational Methods for $C^*$}
\label{sec:az-computation}
%%%%%%%%%%%%%%%%%%%%%%%%%%%%%%%%%%%%%%%%%%%%%%%%%%%%%%%%%%%%%%%%%%%%%%%%%%%%%%%

\subsection{The Computational Challenge}

Computing the fixed point $C^* = \Phi(C^*, \Psi)$ requires iterative methods.
The choice of algorithm depends on the properties of $\Phi$ and the
dimensionality of the problem.

\subsection{Algorithm 1: Fixed-Point Iteration}

The simplest approach is direct iteration:

\begin{tcolorbox}[colback=gray!5!white, colframe=gray!75!black,
    title=Algorithm AZ-C1: Fixed-Point Iteration]
\textbf{Input:} Initial guess $C_0$, context $\Psi$, tolerance $\epsilon$, max iterations $T$

\textbf{Output:} Approximate fixed point $\hat{C}^*$

\begin{enumerate}[nosep]
\item Set $t = 0$
\item \textbf{repeat}
    \begin{enumerate}[nosep]
    \item $C_{t+1} \leftarrow \Phi(C_t, \Psi)$
    \item $\delta_t \leftarrow \|C_{t+1} - C_t\|$
    \item $t \leftarrow t + 1$
    \end{enumerate}
\item \textbf{until} $\delta_t < \epsilon$ or $t > T$
\item \textbf{return} $\hat{C}^* = C_t$
\end{enumerate}

\textbf{Convergence:} Guaranteed if $\|\partial \Phi / \partial C\| < 1$ (contraction).

\textbf{Complexity:} $O(T \cdot n^2)$ where $n = \dim(C)$.
\end{tcolorbox}

\subsection{Algorithm 2: Newton-Raphson Method}

For faster convergence when $\Phi$ is differentiable:

\begin{tcolorbox}[colback=gray!5!white, colframe=gray!75!black,
    title=Algorithm AZ-C2: Newton-Raphson for $C^*$]
\textbf{Input:} Initial guess $C_0$, context $\Psi$, tolerance $\epsilon$

\textbf{Output:} Approximate fixed point $\hat{C}^*$

Define $G(C) = C - \Phi(C, \Psi)$, so $C^*$ solves $G(C^*) = 0$.

\begin{enumerate}[nosep]
\item Set $t = 0$
\item \textbf{repeat}
    \begin{enumerate}[nosep]
    \item Compute Jacobian $J_t = I - \frac{\partial \Phi}{\partial C}\big|_{C_t}$
    \item Solve $J_t \cdot \Delta C = -G(C_t)$
    \item $C_{t+1} \leftarrow C_t + \Delta C$
    \item $t \leftarrow t + 1$
    \end{enumerate}
\item \textbf{until} $\|G(C_t)\| < \epsilon$
\item \textbf{return} $\hat{C}^* = C_t$
\end{enumerate}

\textbf{Convergence:} Quadratic near $C^*$ (much faster than fixed-point).

\textbf{Complexity:} $O(T \cdot n^3)$ due to matrix inversion.
\end{tcolorbox}

\subsection{Algorithm 3: Stochastic Approximation}

When $\Phi$ is observed with noise (empirical settings):

\begin{tcolorbox}[colback=gray!5!white, colframe=gray!75!black,
    title=Algorithm AZ-C3: Robbins-Monro Stochastic Approximation]
\textbf{Input:} Initial guess $C_0$, learning rate sequence $\{\alpha_t\}$

\textbf{Output:} Converged estimate $\hat{C}^*$

\begin{enumerate}[nosep]
\item Set $t = 0$
\item \textbf{repeat}
    \begin{enumerate}[nosep]
    \item Observe noisy $\tilde{\Phi}_t = \Phi(C_t, \Psi_t) + \epsilon_t$
    \item $C_{t+1} \leftarrow C_t + \alpha_t \cdot (\tilde{\Phi}_t - C_t)$
    \item $t \leftarrow t + 1$
    \end{enumerate}
\item \textbf{until} convergence criterion met
\item \textbf{return} $\hat{C}^* = C_t$
\end{enumerate}

\textbf{Learning Rate:} $\alpha_t = a / (b + t)$ with $a > 0$, $b \geq 0$.

\textbf{Convergence:} Almost sure if $\sum \alpha_t = \infty$, $\sum \alpha_t^2 < \infty$.
\end{tcolorbox}

\subsection{Computational Complexity by Context Dimension}

\begin{table}[htbp]
\centering
\caption{Scaling of $C^*$ Computation with $\Psi$-Dimensions}
\label{tab:az-complexity}
\renewcommand{\arraystretch}{1.3}
\begin{tabular}{@{}clll@{}}
\toprule
\textbf{$|\Psi|$} & \textbf{Grid Points} & \textbf{Fixed-Point} & \textbf{Newton} \\
\midrule
2 & $10^2 = 100$ & 0.1 sec & 0.01 sec \\
4 & $10^4 = 10,000$ & 10 sec & 1 sec \\
6 & $10^6$ & 17 min & 2 min \\
8 & $10^8$ & 28 hrs & 3 hrs \\
\bottomrule
\end{tabular}
\end{table}

\textbf{Implication:} For $|\Psi| > 6$, use sparse grids, adaptive methods,
or dimension reduction (PCA on $\Psi$-space).

\subsection{Detecting Multiple Equilibria}

When multiple $C^*$ may exist:

\begin{enumerate}
\item \textbf{Multi-start:} Run Algorithm AZ-C1 from $K$ random initial points
\item \textbf{Homotopy:} Continuously deform $\Psi$ from known solution
\item \textbf{Basin mapping:} Grid search to identify attraction regions
\end{enumerate}

\begin{equation}
\text{Basin}(C^*_j) = \{C_0 : \lim_{t \to \infty} C_t = C^*_j\}
\end{equation}


%%%%%%%%%%%%%%%%%%%%%%%%%%%%%%%%%%%%%%%%%%%%%%%%%%%%%%%%%%%%%%%%%%%%%%%%%%%%%%%
\section{Part 7: Validation and Diagnostics}
\label{sec:az-validation}
%%%%%%%%%%%%%%%%%%%%%%%%%%%%%%%%%%%%%%%%%%%%%%%%%%%%%%%%%%%%%%%%%%%%%%%%%%%%%%%

\subsection{The Validation Challenge}

Any methodological framework must include procedures for validating whether the
model fits data, converges properly, and produces robust results. This section
provides diagnostic tools for EBF.

\subsection{Model Fit Diagnostics}

\begin{tcolorbox}[colback=gray!5!white, colframe=gray!75!black,
    title=Diagnostic Protocol AZ-V1: Model Fit Assessment]
\textbf{Purpose:} Assess how well EBF predictions match observed behavior.

\textbf{Metrics:}
\begin{enumerate}[nosep]
\item \textbf{Root Mean Square Error (RMSE):}
\begin{equation}
\text{RMSE} = \sqrt{\frac{1}{N}\sum_{i=1}^N (U^{obs}_i - U^{pred}_i)^2}
\end{equation}

\item \textbf{Mean Absolute Percentage Error (MAPE):}
\begin{equation}
\text{MAPE} = \frac{100\%}{N}\sum_{i=1}^N \left|\frac{U^{obs}_i - U^{pred}_i}{U^{obs}_i}\right|
\end{equation}

\item \textbf{Pseudo-$R^2$ for discrete choices:}
\begin{equation}
R^2_{McF} = 1 - \frac{\log L_{full}}{\log L_{null}}
\end{equation}
\end{enumerate}

\textbf{Benchmarks:}
\begin{itemize}[nosep]
\item RMSE $< 0.15$: Excellent fit
\item RMSE $\in [0.15, 0.30]$: Acceptable fit
\item RMSE $> 0.30$: Poor fit---investigate model specification
\end{itemize}
\end{tcolorbox}

\subsection{Convergence Diagnostics for $C^*$}

\begin{tcolorbox}[colback=gray!5!white, colframe=gray!75!black,
    title=Diagnostic Protocol AZ-V2: Convergence Checks]
\textbf{Purpose:} Verify that $C^*$ computation has converged properly.

\textbf{Checks:}
\begin{enumerate}[nosep]
\item \textbf{Residual Norm:}
\begin{equation}
\|C_T - \Phi(C_T, \Psi)\| < \epsilon \quad (\epsilon = 10^{-6} \text{ typical})
\end{equation}

\item \textbf{Multi-Start Consistency:}
Run Algorithm AZ-C1 from $K \geq 10$ random starts. Flag if:
\begin{equation}
\max_{j,k} \|C^*_j - C^*_k\| > \delta \quad (\delta = 0.01 \text{ typical})
\end{equation}
This indicates multiple equilibria or non-convergence.

\item \textbf{Iteration Count:}
If $T > T_{max}$ without convergence, algorithm may be stuck. Options:
\begin{itemize}[nosep]
\item Increase $T_{max}$
\item Try Newton-Raphson (AZ-C2)
\item Check for discontinuities in $\Phi$
\end{itemize}

\item \textbf{Eigenvalue Check:}
Compute $\rho(J)$ where $J = \partial\Phi/\partial C|_{C^*}$:
\begin{equation}
\rho(J) < 1 \Rightarrow \text{local stability confirmed}
\end{equation}
\end{enumerate}
\end{tcolorbox}

\subsection{Parameter Sensitivity Analysis}

\begin{tcolorbox}[colback=gray!5!white, colframe=gray!75!black,
    title=Diagnostic Protocol AZ-V3: Sensitivity Analysis]
\textbf{Purpose:} Assess how robust results are to parameter assumptions.

\textbf{Procedure:}
\begin{enumerate}[nosep]
\item \textbf{Local Sensitivity:} Compute elasticities:
\begin{equation}
\eta_{\theta_k} = \frac{\partial \log U}{\partial \log \theta_k}
\end{equation}
High $|\eta_{\theta_k}|$ indicates sensitive parameter.

\item \textbf{Global Sensitivity:} Use Sobol indices:
\begin{equation}
S_k = \frac{\text{Var}(\mathbb{E}[U | \theta_k])}{\text{Var}(U)}
\end{equation}
$\sum_k S_k < 1$ indicates interaction effects; use total indices $S^T_k$.

\item \textbf{Robustness Bounds:} Report results for $\theta \pm 2\sigma_\theta$:
\begin{equation}
U(\theta) \in [U(\theta - 2\sigma), U(\theta + 2\sigma)]
\end{equation}
\end{enumerate}

\textbf{Interpretation:}
\begin{itemize}[nosep]
\item $S_k > 0.3$: Parameter $\theta_k$ is influential---prioritize precise estimation
\item $S_k < 0.05$: Parameter $\theta_k$ is negligible---can use rough estimates
\end{itemize}
\end{tcolorbox}

\subsection{Out-of-Sample Validation}

\begin{tcolorbox}[colback=gray!5!white, colframe=gray!75!black,
    title=Diagnostic Protocol AZ-V4: Cross-Validation]
\textbf{Purpose:} Assess predictive accuracy on unseen data.

\textbf{Procedure:}
\begin{enumerate}[nosep]
\item \textbf{K-Fold Cross-Validation:}
\begin{itemize}[nosep]
\item Split data into $K = 5$ or $K = 10$ folds
\item Train on $K-1$ folds, predict on held-out fold
\item Report mean and std of RMSE across folds
\end{itemize}

\item \textbf{Temporal Validation:}
For panel data, use expanding window:
\begin{equation}
\text{Train: } t \in [1, T_0] \quad \text{Test: } t \in [T_0+1, T_0+h]
\end{equation}
Increase $T_0$ and track prediction accuracy over time.

\item \textbf{Context Generalization:}
Train on context $\Psi_1$, test on $\Psi_2$:
\begin{equation}
\text{RMSE}_{transfer} = \text{RMSE}(\text{predict } \Psi_2 | \text{train } \Psi_1)
\end{equation}
High transfer RMSE indicates context-specific parameters need re-estimation.
\end{enumerate}
\end{tcolorbox}

\subsection{Diagnostic Summary Table}

\begin{table}[htbp]
\centering
\caption{Validation Diagnostic Checklist}
\label{tab:az-diagnostics}
\renewcommand{\arraystretch}{1.3}
\begin{tabular}{@{}llll@{}}
\toprule
\textbf{Diagnostic} & \textbf{Metric} & \textbf{Threshold} & \textbf{Action if Fail} \\
\midrule
\multicolumn{4}{@{}l}{\textit{Model Fit}} \\
In-sample RMSE & $\sqrt{\text{MSE}}$ & $< 0.30$ & Re-specify model \\
Pseudo-$R^2$ & $R^2_{McF}$ & $> 0.20$ & Add covariates \\
\addlinespace
\multicolumn{4}{@{}l}{\textit{Convergence}} \\
Residual norm & $\|C - \Phi(C)\|$ & $< 10^{-6}$ & Increase iterations \\
Multi-start agreement & $\max\|C^*_j - C^*_k\|$ & $< 0.01$ & Multiple equilibria \\
Spectral radius & $\rho(J)$ & $< 1$ & Unstable $C^*$ \\
\addlinespace
\multicolumn{4}{@{}l}{\textit{Robustness}} \\
Sobol first-order & $S_k$ & Report all & Prioritize high-$S_k$ \\
Parameter bounds & $U(\theta \pm 2\sigma)$ & Report range & Widen CI if large \\
\addlinespace
\multicolumn{4}{@{}l}{\textit{Prediction}} \\
Cross-validation RMSE & CV-RMSE & $< 1.2 \times$ in-sample & Overfitting \\
Temporal stability & $\Delta$RMSE over time & $< 0.05$/period & Model drift \\
\bottomrule
\end{tabular}
\end{table}

\subsection{Red Flags and Common Problems}

\begin{tcolorbox}[colback=red!5!white, colframe=red!75!black,
    title=Warning Signs in EBF Estimation]

\textbf{Red Flag 1: Non-Convergence}
\begin{itemize}[nosep]
\item \textit{Symptom:} Algorithm oscillates, residual doesn't decrease
\item \textit{Cause:} $\Phi$ may be non-contractive or discontinuous
\item \textit{Fix:} Use damped iteration: $C_{t+1} = \alpha \Phi(C_t) + (1-\alpha)C_t$
\end{itemize}

\textbf{Red Flag 2: Implausible Parameters}
\begin{itemize}[nosep]
\item \textit{Symptom:} $\hat{\gamma} > 1$ or $\hat{\lambda} < 0$
\item \textit{Cause:} Identification failure, multicollinearity
\item \textit{Fix:} Add constraints, check data quality, use regularization
\end{itemize}

\textbf{Red Flag 3: Context Collapse}
\begin{itemize}[nosep]
\item \textit{Symptom:} $C^*(\Psi_1) \approx C^*(\Psi_2)$ for very different $\Psi$
\item \textit{Cause:} Context dimensions not measured properly
\item \textit{Fix:} Refine $\Psi$ measurement, add interaction terms
\end{itemize}

\textbf{Red Flag 4: Overfitting}
\begin{itemize}[nosep]
\item \textit{Symptom:} CV-RMSE $\gg$ in-sample RMSE
\item \textit{Cause:} Too many parameters relative to data
\item \textit{Fix:} Regularization, dimension reduction, simpler model
\end{itemize}
\end{tcolorbox}


%%%%%%%%%%%%%%%%%%%%%%%%%%%%%%%%%%%%%%%%%%%%%%%%%%%%%%%%%%%%%%%%%%%%%%%%%%%%%%%
\section{Part 8: Software Implementation}
\label{sec:az-software}
%%%%%%%%%%%%%%%%%%%%%%%%%%%%%%%%%%%%%%%%%%%%%%%%%%%%%%%%%%%%%%%%%%%%%%%%%%%%%%%

\subsection{Implementation Overview}

This section provides pseudocode implementations for the core EBF methods.
These can be translated to R, Python, Julia, or other languages.

\subsection{Core Data Structures}

\begin{tcolorbox}[colback=gray!5!white, colframe=gray!75!black,
    title=Data Structure: EBF Parameters]
\begin{verbatim}
STRUCT BCM_Parameters:
    # FEPSDE dimension weights
    omega[6]          # F, E, P, S, D, Ec weights

    # Complementarity matrix (symmetric)
    gamma[6][6]       # gamma[d1][d2] for dimensions
    gamma_cross[L][L] # cross-level complementarity

    # Loss aversion by dimension
    lambda[6]         # lambda_F, lambda_E, etc.

    # Time preferences
    delta             # discount factor
    beta              # present bias (quasi-hyperbolic)

    # AWX-WAX parameters
    awareness         # A in [0,1]
    threshold         # theta action threshold
    phi               # thinking style in [0,1]

    # Context
    psi               # context vector
\end{verbatim}
\end{tcolorbox}

\subsection{Algorithm: Compute $C^*$ (Python-style Pseudocode)}

\begin{tcolorbox}[colback=gray!5!white, colframe=gray!75!black,
    title=Implementation AZ-S1: Fixed-Point Iteration for $C^*$]
\begin{verbatim}
def compute_Cstar(psi, params, tol=1e-6, max_iter=1000):
    """
    Compute self-consistent reference structure C*.

    Args:
        psi: Context vector (dict or array)
        params: BCM_Parameters object
        tol: Convergence tolerance
        max_iter: Maximum iterations

    Returns:
        C_star: Converged reference structure
        converged: Boolean convergence flag
    """
    # Initialize C at neutral point
    C = initialize_reference(psi)

    for t in range(max_iter):
        # Apply update operator Phi
        C_new = Phi(C, psi, params)

        # Check convergence
        delta = norm(C_new - C)
        if delta < tol:
            return C_new, True

        # Damped update for stability
        alpha = 0.7  # damping factor
        C = alpha * C_new + (1 - alpha) * C

    return C, False  # Did not converge

def Phi(C, psi, params):
    """
    Update operator: maps current C to implied C.
    """
    # Compute expected utility given C as reference
    U_expected = compute_utility(C, psi, params)

    # Reference adjusts toward expected
    C_new = weighted_average(C, U_expected,
                             weight=params.adjustment_speed)
    return C_new
\end{verbatim}
\end{tcolorbox}

\subsection{Algorithm: Utility Computation with Complementarity}

\begin{tcolorbox}[colback=gray!5!white, colframe=gray!75!black,
    title=Implementation AZ-S2: EBF Utility Function]
\begin{verbatim}
def compute_utility(outcomes, psi, params):
    """
    Compute EBF utility with complementarity.

    Args:
        outcomes: Dict with FEPSDE outcomes {F, E, P, S, D, Ec}
        psi: Context vector
        params: BCM_Parameters

    Returns:
        U: Total utility (scalar)
    """
    dims = ['F', 'E', 'P', 'S', 'D', 'Ec']

    # 1. Compute dimension-specific utilities
    U_dim = {}
    for d in dims:
        # Reference-dependent value function
        x = outcomes[d]
        r = params.reference[d]

        if x >= r:
            U_dim[d] = (x - r) ** params.alpha
        else:
            U_dim[d] = -params.lambda[d] * (r - x) ** params.beta

    # 2. Additive component
    U_additive = sum(params.omega[d] * U_dim[d] for d in dims)

    # 3. Complementarity component
    U_complement = 0
    for i, d1 in enumerate(dims):
        for j, d2 in enumerate(dims):
            if i < j:
                U_complement += params.gamma[i][j] * U_dim[d1] * U_dim[d2]

    # 4. Combine based on thinking style phi
    # phi=0: purely additive, phi=1: full complementarity
    U_total = (1 - params.phi) * U_additive + params.phi * (U_additive + U_complement)

    return U_total
\end{verbatim}
\end{tcolorbox}

\subsection{Algorithm: Heterogeneous Population Simulation}

\begin{tcolorbox}[colback=gray!5!white, colframe=gray!75!black,
    title=Implementation AZ-S3: Population Behavior Simulation]
\begin{verbatim}
def simulate_population(N, psi, param_distributions, intervention=None):
    """
    Simulate behavior for heterogeneous population.

    Args:
        N: Population size
        psi: Context vector
        param_distributions: Dict of parameter distributions
        intervention: Optional intervention to apply

    Returns:
        results: DataFrame with individual outcomes
    """
    results = []

    for i in range(N):
        # Draw individual parameters from distributions
        params_i = draw_parameters(param_distributions)

        # Compute C* for this individual
        C_star_i, _ = compute_Cstar(psi, params_i)

        # Compute potential utility
        U_pot = compute_utility(outcomes, psi, params_i)

        # Apply awareness filter (AWX)
        U_eff = params_i.awareness * U_pot

        # Compute willingness (WAX)
        WAX = compute_WAX(U_eff, params_i)

        # Action decision
        action = WAX >= params_i.threshold

        # Store results
        results.append({
            'id': i,
            'params': params_i,
            'C_star': C_star_i,
            'U_pot': U_pot,
            'U_eff': U_eff,
            'WAX': WAX,
            'action': action
        })

    return pd.DataFrame(results)

def draw_parameters(distributions):
    """Draw one set of parameters from distributions."""
    params = BCM_Parameters()

    # Example: threshold from LogNormal
    params.threshold = np.random.lognormal(
        distributions['theta_mu'],
        distributions['theta_sigma']
    )

    # Example: awareness from Beta
    params.awareness = np.random.beta(
        distributions['A_alpha'],
        distributions['A_beta']
    )

    # ... other parameters ...
    return params
\end{verbatim}
\end{tcolorbox}

\subsection{Algorithm: Learning Dynamics Simulation}

\begin{tcolorbox}[colback=gray!5!white, colframe=gray!75!black,
    title=Implementation AZ-S4: Adaptive Learning Simulation]
\begin{verbatim}
def simulate_learning(T, psi_sequence, params, noise_sd=0.1):
    """
    Simulate learning dynamics for C* adjustment.

    Args:
        T: Number of time periods
        psi_sequence: List of context vectors over time
        params: BCM_Parameters (initial)
        noise_sd: Standard deviation of observation noise

    Returns:
        C_path: Array of C values over time
        converged_at: Time of convergence (or None)
    """
    C_path = np.zeros((T, len(params.reference)))
    C = params.reference.copy()  # Initial reference

    # Learning rate schedule: alpha_t = a / (b + t)
    a, b = 1.0, 10.0

    for t in range(T):
        psi_t = psi_sequence[t]

        # Observe noisy feedback
        C_target = Phi(C, psi_t, params)
        C_observed = C_target + np.random.normal(0, noise_sd, len(C))

        # Adaptive update (Robbins-Monro)
        alpha_t = a / (b + t)
        C = C + alpha_t * (C_observed - C)

        C_path[t] = C

        # Check convergence
        if t > 10:
            recent_change = np.std(C_path[t-10:t], axis=0).max()
            if recent_change < 0.001:
                return C_path[:t+1], t

    return C_path, None
\end{verbatim}
\end{tcolorbox}

\subsection{Algorithm: Parameter Estimation via Maximum Likelihood}

\begin{tcolorbox}[colback=gray!5!white, colframe=gray!75!black,
    title=Implementation AZ-S5: Maximum Likelihood Estimation]
\begin{verbatim}
def estimate_parameters(data, initial_params, method='BFGS'):
    """
    Estimate EBF parameters via MLE.

    Args:
        data: DataFrame with observed choices and contexts
        initial_params: Starting parameter values
        method: Optimization method

    Returns:
        params_hat: Estimated parameters
        se: Standard errors
        diagnostics: Fit statistics
    """
    def neg_log_likelihood(theta):
        """Negative log-likelihood to minimize."""
        params = vector_to_params(theta)
        ll = 0

        for _, row in data.iterrows():
            # Compute choice probability
            WAX = compute_WAX(row['outcomes'], row['psi'], params)
            p_action = sigmoid(WAX - params.threshold)

            # Add to log-likelihood
            if row['action']:
                ll += np.log(p_action + 1e-10)
            else:
                ll += np.log(1 - p_action + 1e-10)

        return -ll

    # Optimize
    theta_init = params_to_vector(initial_params)
    result = minimize(neg_log_likelihood, theta_init,
                      method=method, options={'disp': True})

    # Extract results
    params_hat = vector_to_params(result.x)

    # Compute standard errors via Hessian
    hessian = compute_hessian(neg_log_likelihood, result.x)
    se = np.sqrt(np.diag(np.linalg.inv(hessian)))

    # Diagnostics
    diagnostics = {
        'log_likelihood': -result.fun,
        'AIC': 2 * len(theta_init) + 2 * result.fun,
        'BIC': np.log(len(data)) * len(theta_init) + 2 * result.fun,
        'converged': result.success
    }

    return params_hat, se, diagnostics
\end{verbatim}
\end{tcolorbox}

\subsection{Software Implementation Notes}

\begin{tcolorbox}[colback=green!5!white, colframe=green!75!black,
    title=Implementation Guidelines]
\begin{enumerate}
\item \textbf{Language Choice:}
\begin{itemize}[nosep]
\item Python: Best for prototyping, integration with LLM APIs
\item R: Best for statistical analysis, existing econometric packages
\item Julia: Best for computational speed, large-scale simulations
\end{itemize}

\item \textbf{Numerical Stability:}
\begin{itemize}[nosep]
\item Use $\log$-likelihoods, not raw probabilities
\item Clip extreme values to avoid overflow/underflow
\item Use damped iteration for $C^*$ computation
\end{itemize}

\item \textbf{Testing:}
\begin{itemize}[nosep]
\item Unit test each component separately
\item Verify $C^* = \Phi(C^*, \Psi)$ holds at convergence
\item Check parameter recovery on simulated data
\end{itemize}

\item \textbf{Performance:}
\begin{itemize}[nosep]
\item Vectorize operations where possible
\item Use sparse matrices for large $\Gamma$
\item Parallelize population simulations
\end{itemize}
\end{enumerate}
\end{tcolorbox}


%%%%%%%%%%%%%%%%%%%%%%%%%%%%%%%%%%%%%%%%%%%%%%%%%%%%%%%%%%%%%%%%%%%%%%%%%%%%%%%
\section{Part 9: Bayesian Prior Specification}
\label{sec:az-priors}
%%%%%%%%%%%%%%%%%%%%%%%%%%%%%%%%%%%%%%%%%%%%%%%%%%%%%%%%%%%%%%%%%%%%%%%%%%%%%%%

\subsection{The Prior Specification Challenge}

Bayesian estimation requires specifying prior distributions for all parameters.
Poor priors can dominate results, especially with limited data. This section
provides guidance for EBF prior specification.

\subsection{Prior Types and Recommendations}

\begin{tcolorbox}[colback=gray!5!white, colframe=gray!75!black,
    title=Prior Specification: Core Parameters]
\textbf{1. Complementarity $\gamma_{dd'}$:}
\begin{equation}
\gamma_{dd'} \sim \text{Normal}(0.3, 0.2^2) \quad \text{truncated to } [-1, 1]
\end{equation}
\begin{itemize}[nosep]
\item Rationale: Empirical estimates cluster around 0.2--0.5 for positive complements
\item Truncation ensures interpretable range
\item Weakly informative: allows data to dominate
\end{itemize}

\textbf{2. Loss Aversion $\lambda$:}
\begin{equation}
\lambda \sim \text{LogNormal}(\ln(2.25), 0.4^2)
\end{equation}
\begin{itemize}[nosep]
\item Mode at $\approx 2.25$ (Tversky-Kahneman baseline)
\item Log-normal ensures $\lambda > 0$
\item Right-skewed to allow for high loss aversion
\end{itemize}

\textbf{3. Dimension Weights $\omega_d$:}
\begin{equation}
\boldsymbol{\omega} \sim \text{Dirichlet}(\alpha_1, \ldots, \alpha_6) \quad \text{with } \alpha_d = 2
\end{equation}
\begin{itemize}[nosep]
\item Ensures $\sum_d \omega_d = 1$
\item $\alpha_d = 2$ is weakly informative (mild regularization)
\item Can set $\alpha_d$ differently if prior knowledge exists
\end{itemize}

\textbf{4. Discount Factor $\delta$:}
\begin{equation}
\delta \sim \text{Beta}(9, 1) \quad \text{mapped to } [0.5, 1]
\end{equation}
\begin{itemize}[nosep]
\item Prior mode around 0.9 (annual discount rate 10\%)
\item Constrained to economically plausible range
\end{itemize}
\end{tcolorbox}

\begin{tcolorbox}[colback=gray!5!white, colframe=gray!75!black,
    title=Prior Specification: AWX-WAX Parameters]
\textbf{5. Awareness $A$:}
\begin{equation}
A \sim \text{Beta}(2, 3)
\end{equation}
\begin{itemize}[nosep]
\item Mode at $\approx 0.33$, mean $\approx 0.4$
\item Reflects that most people don't fully attend to all utility dimensions
\item Bounded on $[0, 1]$ by definition
\end{itemize}

\textbf{6. Action Threshold $\theta$:}
\begin{equation}
\theta \sim \text{LogNormal}(-0.5, 0.5^2)
\end{equation}
\begin{itemize}[nosep]
\item Median $\approx 0.6$
\item Log-normal captures right-skew (some have very high thresholds)
\item Consistent with intention-behavior gap literature
\end{itemize}

\textbf{7. Thinking Style $\varphi$:}
\begin{equation}
\varphi \sim \text{Beta}(2, 2)
\end{equation}
\begin{itemize}[nosep]
\item Symmetric around 0.5
\item Allows full range from additive ($\varphi=0$) to complementary ($\varphi=1$)
\item Mildly regularizing
\end{itemize}
\end{tcolorbox}

\subsection{Prior Elicitation from LLM}

\begin{tcolorbox}[colback=blue!5!white, colframe=blue!75!black,
    title=LLM-Based Prior Elicitation Protocol]
When domain expertise is unavailable, LLMs can help specify informative priors:

\textbf{Step 1: Describe the Parameter}
\begin{verbatim}
"In behavioral economics, loss aversion (lambda) measures
how much more painful losses are compared to equivalent
gains. A lambda of 2 means losses hurt twice as much."
\end{verbatim}

\textbf{Step 2: Ask for Distribution}
\begin{verbatim}
"Based on the behavioral economics literature, what is a
reasonable prior distribution for lambda? Provide:
- A central estimate (mode or median)
- A plausible range (95% interval)
- The shape (symmetric, skewed, bounded?)"
\end{verbatim}

\textbf{Step 3: Convert to Parametric Form}
\begin{verbatim}
"Given mode=2.25, 95% CI=[1.5, 4.0], and right-skew,
suggest a parametric distribution (e.g., LogNormal,
Gamma) with specific parameters."
\end{verbatim}

\textbf{Step 4: Validate}
\begin{itemize}[nosep]
\item Check implied prior predictive distribution
\item Compare to known empirical benchmarks
\item Run sensitivity analysis (see AZ-V3)
\end{itemize}
\end{tcolorbox}

\subsection{Hierarchical Priors for Heterogeneity}

When estimating heterogeneous populations:

\begin{equation}
\theta_i \sim \text{Normal}(\mu_\theta, \sigma_\theta^2)
\end{equation}
\begin{equation}
\mu_\theta \sim \text{Normal}(0.5, 0.3^2) \quad \sigma_\theta \sim \text{HalfNormal}(0.3)
\end{equation}

\begin{tcolorbox}[colback=yellow!5!white, colframe=yellow!75!black,
    title=Hierarchical Prior Strategy]
\textbf{Level 1 (Individual):} $\theta_i | \mu_\theta, \sigma_\theta$

\textbf{Level 2 (Population):} $\mu_\theta, \sigma_\theta | \text{hyperpriors}$

\textbf{Benefits:}
\begin{itemize}[nosep]
\item Partial pooling: Individual estimates shrink toward population mean
\item Uncertainty propagation: Population uncertainty affects individual estimates
\item Identifiability: Can estimate even with few observations per individual
\end{itemize}

\textbf{Hyperprior Guidance:}
\begin{itemize}[nosep]
\item Use weakly informative hyperpriors
\item $\sigma_\theta$: HalfNormal or HalfCauchy(0, $s$) where $s$ is plausible SD
\item $\mu_\theta$: Centered on literature estimate with moderate variance
\end{itemize}
\end{tcolorbox}

\subsection{Prior Sensitivity Analysis}

\begin{table}[htbp]
\centering
\caption{Prior Sensitivity Checklist}
\label{tab:az-prior-sensitivity}
\renewcommand{\arraystretch}{1.3}
\begin{tabular}{@{}llp{5cm}@{}}
\toprule
\textbf{Check} & \textbf{Method} & \textbf{Interpretation} \\
\midrule
Prior predictive & Simulate from prior & Does prior imply plausible predictions? \\
Prior-posterior overlap & Compare distributions & High overlap = data uninformative \\
Different priors & Re-estimate with alternatives & Results robust to prior choice? \\
Effective sample size & Compute ESS & ESS $<$ nominal $N$ if prior dominates \\
\bottomrule
\end{tabular}
\end{table}

\paragraph{Prior Predictive Check.}
Before seeing data, simulate outcomes from the prior:
\begin{enumerate}
\item Draw $\theta^{(s)} \sim p(\theta)$ for $s = 1, \ldots, S$
\item Simulate $y^{(s)} \sim p(y | \theta^{(s)})$
\item Check if $y^{(s)}$ values are plausible given domain knowledge
\end{enumerate}

If prior predictive includes implausible values (e.g., negative utility shares,
$\lambda < 1$), revise the prior.

\subsection{Summary: Recommended Prior Table}

\begin{table}[htbp]
\centering
\caption{Recommended Priors for EBF Parameters}
\label{tab:az-recommended-priors}
\renewcommand{\arraystretch}{1.3}
\small
\begin{tabular}{@{}llll@{}}
\toprule
\textbf{Parameter} & \textbf{Distribution} & \textbf{Parameters} & \textbf{Source} \\
\midrule
$\gamma_{dd'}$ & Truncated Normal & $\mu=0.3$, $\sigma=0.2$, $[-1,1]$ & Literature synthesis \\
$\lambda$ & LogNormal & $\mu=\ln(2.25)$, $\sigma=0.4$ & Tversky \& Kahneman \\
$\boldsymbol{\omega}$ & Dirichlet & $\boldsymbol{\alpha} = (2,2,2,2,2,2)$ & Weakly informative \\
$\delta$ & Beta (scaled) & $\alpha=9$, $\beta=1$ on $[0.5,1]$ & Frederick et al. \\
$A$ & Beta & $\alpha=2$, $\beta=3$ & Salience literature \\
$\theta$ & LogNormal & $\mu=-0.5$, $\sigma=0.5$ & Intention-behavior gap \\
$\varphi$ & Beta & $\alpha=2$, $\beta=2$ & CRT distribution \\
\bottomrule
\end{tabular}
\end{table}

\begin{tcolorbox}[colback=green!5!white, colframe=green!75!black,
    title=Best Practices for Prior Specification]
\begin{enumerate}
\item \textbf{Start weakly informative:} Let data speak unless strong prior knowledge exists
\item \textbf{Respect constraints:} Use distributions that respect parameter bounds
\item \textbf{Document choices:} Record rationale for each prior
\item \textbf{Run sensitivity:} Always check how results change with alternative priors
\item \textbf{Update iteratively:} Refine priors as evidence accumulates across studies
\end{enumerate}
\end{tcolorbox}


%%%%%%%%%%%%%%%%%%%%%%%%%%%%%%%%%%%%%%%%%%%%%%%%%%%%%%%%%%%%%%%%%%%%%%%%%%%%%%%
\section{Integration with Other Appendices}
\label{sec:az-integration}
%%%%%%%%%%%%%%%%%%%%%%%%%%%%%%%%%%%%%%%%%%%%%%%%%%%%%%%%%%%%%%%%%%%%%%%%%%%%%%%

\subsection{Connection to FORMAL-PROOF (D)}

\begin{tcolorbox}[colback=yellow!5!white, colframe=yellow!75!black,
    title=Integration: METHOD-CONSTRUCT $\leftrightarrow$ FORMAL-PROOF (D)]
\textbf{What this appendix provides:}
\begin{itemize}[nosep]
\item Methodological justification for formal axioms
\item Learning dynamics for parameter convergence
\item Identification conditions for empirical testing
\end{itemize}

\textbf{What it requires:}
\begin{itemize}[nosep]
\item Complete proofs of $C^*$ existence and uniqueness (D)
\item Stability conditions for fixed-point (D)
\item Mathematical regularity assumptions (D)
\end{itemize}
\end{tcolorbox}

\subsection{Connection to CORE-WHERE (BBB)}

\begin{tcolorbox}[colback=yellow!5!white, colframe=yellow!75!black,
    title=Integration: METHOD-CONSTRUCT $\leftrightarrow$ CORE-WHERE (BBB)]
\textbf{What this appendix provides:}
\begin{itemize}[nosep]
\item Identification conditions for parameter estimation
\item Aggregation rules for population parameters
\item Learning convergence for calibration
\end{itemize}

\textbf{What it requires:}
\begin{itemize}[nosep]
\item Epistemic status hierarchy (EMP, THR, LLM, ILL, HYP)
\item Prior specification methodology
\item Validation protocols
\end{itemize}
\end{tcolorbox}

\subsection{Connection to METHOD-LLMMC (AN)}

\begin{tcolorbox}[colback=yellow!5!white, colframe=yellow!75!black,
    title=Integration: METHOD-CONSTRUCT $\leftrightarrow$ METHOD-LLMMC (AN)]
\textbf{What this appendix provides:}
\begin{itemize}[nosep]
\item Theoretical grounding for LLM as ``synthetic sample''
\item Convergence conditions for Monte Carlo estimation
\item Identification of what LLM elicitation can and cannot identify
\end{itemize}

\textbf{What it requires:}
\begin{itemize}[nosep]
\item Practical LLM prompting methodology
\item Variance estimation procedures
\item Calibration against empirical benchmarks
\end{itemize}
\end{tcolorbox}

\subsection{Connection to CORE-AWARE (AU) and CORE-READY (AV)}

The 9C framework completes the path from utility to action. METHOD-CONSTRUCT
provides the methodological foundations for this chain:

\begin{tcolorbox}[colback=yellow!5!white, colframe=yellow!75!black,
    title=Integration: METHOD-CONSTRUCT $\leftrightarrow$ CORE-AWARE/READY (AU, AV)]
\textbf{What this appendix provides:}
\begin{itemize}[nosep]
\item Microfoundation for the Awareness filter $A(\cdot)$
\item Learning dynamics for salience updating
\item Identification of $\theta$ (action threshold) from behavior
\item Heterogeneity structure for $\varphi$ (thinking style)
\end{itemize}

\textbf{What it requires:}
\begin{itemize}[nosep]
\item Complete AWX axiom system (AU): $U_{eff} = A \cdot U_{pot}$
\item Complete WAX axiom system (AV): Action $\Leftrightarrow WAX \geq \theta$
\item Thinking style specification: $\varphi \in [0,1]$
\end{itemize}
\end{tcolorbox}

\paragraph{The AWX$\to$WAX Pipeline.}
The full behavioral chain requires methodological grounding at each step:

\begin{enumerate}
\item \textbf{Potential Utility} $U_{pot}$: From CORE-WHAT (C), identified via Part 5
\item \textbf{Awareness Filter} $A(t^*)$: From CORE-AWARE (AU), learned via Part 2
\item \textbf{Effective Utility} $U_{eff} = A \cdot U_{pot}$: Multiplicative, testable
\item \textbf{Willingness} $WAX$: From CORE-READY (AV), aggregated via Part 4
\item \textbf{Threshold} $\theta$: Heterogeneous, distributed via AZ-H1
\item \textbf{Action}: Binary outcome, $WAX \geq \theta$
\end{enumerate}

\paragraph{Methodological Implications.}

\begin{itemize}
\item \textbf{Awareness Dynamics:} Salience $\sigma$ evolves via Bayesian updating
(AZ-L1), with decay rate $\tau$ estimated from attention data.

\item \textbf{Threshold Identification:} $\theta_i$ is identified from the
discontinuity at $WAX = \theta$ in behavior probability.

\item \textbf{Thinking Style:} The parameter $\varphi$ (additive vs. complementary
thinking) is identified from interaction effects in utility responses.
\end{itemize}

\begin{equation}
\boxed{
P(\text{Action}) = \int \mathbf{1}[WAX_i(A_i \cdot U_{pot}, \varphi_i) \geq \theta_i] \, dF(\Theta_i)
}
\end{equation}

\subsection{Connection to METHOD-EVAL (R): Validation Framework}

\begin{tcolorbox}[colback=yellow!5!white, colframe=yellow!75!black,
    title=Integration: METHOD-CONSTRUCT $\leftrightarrow$ METHOD-EVAL (R)]
\textbf{What this appendix provides:}
\begin{itemize}[nosep]
\item Diagnostic protocols (AZ-V1 through AZ-V4) for model evaluation
\item Convergence criteria for $C^*$ computation
\item Sensitivity analysis framework (Sobol indices)
\item Red flags and common problems checklist
\end{itemize}

\textbf{What it requires:}
\begin{itemize}[nosep]
\item Evaluation criteria for EBF predictions (R)
\item Benchmark datasets for validation (R)
\item Success metrics for behavioral interventions (R)
\end{itemize}
\end{tcolorbox}

\paragraph{Validation Pipeline.}
The complete validation workflow integrates both appendices:
\begin{enumerate}
\item \textbf{Estimate} parameters using METHOD-LLMMC (AN) or traditional methods
\item \textbf{Diagnose} using AZ protocols: fit, convergence, sensitivity
\item \textbf{Evaluate} against R criteria: predictive accuracy, calibration
\item \textbf{Iterate} if diagnostics fail
\end{enumerate}

\subsection{Connection to METHOD-OPS (E): Operationalization}

\begin{tcolorbox}[colback=yellow!5!white, colframe=yellow!75!black,
    title=Integration: METHOD-CONSTRUCT $\leftrightarrow$ METHOD-OPS (E)]
\textbf{What this appendix provides:}
\begin{itemize}[nosep]
\item Theoretical constructs requiring measurement (FEPSDE, $\gamma$, $\Psi$)
\item Identification conditions for what can be measured
\item Heterogeneity structure guiding sample stratification
\end{itemize}

\textbf{What it requires:}
\begin{itemize}[nosep]
\item Survey instruments for FEPSDE dimensions (E)
\item Context ($\Psi$) measurement protocols (E)
\item Scale validation procedures (E)
\end{itemize}
\end{tcolorbox}

\paragraph{From Theory to Measurement.}
The operationalization chain:
\begin{equation}
\text{AZ: Constructs} \xrightarrow{\text{E: Instruments}} \text{AN: Estimation} \xrightarrow{\text{R: Validation}} \text{BBB: Calibrated Values}
\end{equation}

\subsection{Connection to METHOD-SRL (AL): Self-Reinforcement Learning}

\begin{tcolorbox}[colback=yellow!5!white, colframe=yellow!75!black,
    title=Integration: METHOD-CONSTRUCT $\leftrightarrow$ METHOD-SRL (AL)]
\textbf{What this appendix provides:}
\begin{itemize}[nosep]
\item Learning dynamics framework (AZ-L1, AZ-L2)
\item Convergence conditions for adaptive learning
\item Level-$k$ thinking structure
\end{itemize}

\textbf{What it requires:}
\begin{itemize}[nosep]
\item Self-reinforcement mechanisms in behavior change (AL)
\item Habit formation dynamics (AL)
\item Feedback loop modeling (AL)
\end{itemize}
\end{tcolorbox}

\paragraph{Learning Integration.}
AZ provides the general Bayesian learning framework; AL specializes it for
self-reinforcing behaviors where past actions affect future preferences:
\begin{equation}
U_{t+1}(\text{Action}) = U_t(\text{Action}) + \beta \cdot \mathbf{1}[\text{Action}_t = \text{Action}]
\end{equation}

\subsection{METHOD Appendix Ecosystem}

\begin{table}[htbp]
\centering
\caption{METHOD Appendix Integration Map}
\label{tab:az-method-ecosystem}
\renewcommand{\arraystretch}{1.3}
\begin{tabular}{@{}lllp{4cm}@{}}
\toprule
\textbf{From} & \textbf{To} & \textbf{Flow} & \textbf{Content} \\
\midrule
AZ & AN & Theory $\to$ Estimation & Identification conditions, algorithms \\
AZ & R & Theory $\to$ Validation & Diagnostic protocols, criteria \\
AZ & E & Theory $\to$ Measurement & Construct definitions, scales \\
AZ & AL & Theory $\to$ Learning & Learning axioms, convergence \\
\addlinespace
AN & AZ & Estimation $\to$ Theory & Empirical parameter values \\
R & AZ & Validation $\to$ Theory & Model fit feedback \\
E & AZ & Measurement $\to$ Theory & Operationalized constructs \\
AL & AZ & Learning $\to$ Theory & Reinforcement dynamics \\
\bottomrule
\end{tabular}
\end{table}

\paragraph{Practitioner Workflow.}
For a complete EBF application:
\begin{enumerate}
\item \textbf{Ground in AZ:} Understand methodological foundations
\item \textbf{Operationalize via E:} Create measurement instruments
\item \textbf{Estimate via AN:} Use LLM-MC or traditional methods
\item \textbf{Validate via R:} Check predictions against benchmarks
\item \textbf{Model dynamics via AL:} If behavior change is the goal
\item \textbf{Calibrate to BBB:} Store validated parameter values
\end{enumerate}


\subsection{Connection to Behavioral Change Journey (Chapters 13--16)}

The Behavioral Change Journey extends EBF from static decisions to dynamic
change processes. METHOD-CONSTRUCT provides the methodological foundations:

\begin{tcolorbox}[colback=yellow!5!white, colframe=yellow!75!black,
    title=Integration: METHOD-CONSTRUCT $\leftrightarrow$ Behavioral Change]
\textbf{What this appendix provides:}
\begin{itemize}[nosep]
\item Dynamic adjustment framework for stage transitions (Part 3)
\item Heterogeneous segment modeling (Part 4, AZ-H1)
\item Identification of segment membership from behavior (Part 5)
\item Learning dynamics for habit formation (Part 2)
\end{itemize}

\textbf{What it requires:}
\begin{itemize}[nosep]
\item Stage definitions from Ch. 13 (Pre-contemplation $\to$ Action $\to$ Maintenance)
\item Segment profiles from Ch. 14 (Resisters, Deliberators, Actives)
\item WEC synthesis from Ch. 15 (Willingness $\times$ Journey $\times$ Segment)
\item Probability of Change output from Ch. 16
\end{itemize}
\end{tcolorbox}

\paragraph{Stage Transition Dynamics.}
Stage transitions follow the adjustment dynamics from AZ-D1:
\begin{equation}
P(\text{Stage}_{t+1} = s' | \text{Stage}_t = s) = f(WAX_i, \theta_i, \Delta C^*)
\end{equation}

The transition probability depends on:
\begin{itemize}[nosep]
\item $WAX_i$: Current willingness (from CORE-READY)
\item $\theta_i$: Individual threshold (heterogeneous, from AZ-H1)
\item $\Delta C^*$: Reference shift induced by intervention
\end{itemize}

\paragraph{Segment-Specific Parameters.}
Different behavioral segments have different parameter distributions:

\begin{table}[htbp]
\centering
\caption{Parameter Profiles by Behavioral Segment}
\label{tab:az-segments}
\renewcommand{\arraystretch}{1.2}
\begin{tabular}{@{}lccc@{}}
\toprule
\textbf{Parameter} & \textbf{Resisters} & \textbf{Deliberators} & \textbf{Actives} \\
\midrule
$\theta$ (threshold) & High (0.7) & Medium (0.5) & Low (0.3) \\
$\alpha$ (learning rate) & Low (0.03) & Medium (0.08) & High (0.15) \\
$A$ (awareness) & Low (0.25) & Medium (0.45) & High (0.70) \\
$\varphi$ (complementary) & Low (0.2) & Medium (0.5) & High (0.7) \\
\bottomrule
\end{tabular}
\end{table}

\paragraph{Methodological Implication.}
Interventions should be segment-specific:
\begin{itemize}
\item \textbf{Resisters:} Lower $\theta$ through friction reduction
\item \textbf{Deliberators:} Increase $A$ through salience interventions
\item \textbf{Actives:} Maintain behavior through habit reinforcement ($\alpha \uparrow$)
\end{itemize}


% =============================================================================
%                    PART 3: APPLICATION
% =============================================================================

%%%%%%%%%%%%%%%%%%%%%%%%%%%%%%%%%%%%%%%%%%%%%%%%%%%%%%%%%%%%%%%%%%%%%%%%%%%%%%%
\section{Worked Example}
\label{sec:az-example}
%%%%%%%%%%%%%%%%%%%%%%%%%%%%%%%%%%%%%%%%%%%%%%%%%%%%%%%%%%%%%%%%%%%%%%%%%%%%%%%

\subsection{Example: Identifying Complementarity in Financial-Emotional Dimensions}

\paragraph{Setup.}
Consider estimating $\gamma_{FE}$ (Financial-Emotional complementarity) using
within-person variation. A sample of $N = 500$ individuals reports well-being
under two conditions: (1) financial gain alone, (2) financial gain with emotional
support.

\paragraph{Application.}

\textbf{Step 1: Specification}
\begin{equation}
U_i = \omega_F \cdot U_F + \omega_E \cdot U_E + \gamma_{FE} \cdot U_F \cdot U_E + \epsilon_i
\end{equation}

\textbf{Step 2: Identification Strategy}

Define treatment $T_i = 1$ if emotional support present:
\begin{align}
U_i(T=0) &= \omega_F \cdot U_F + \epsilon_i \\
U_i(T=1) &= \omega_F \cdot U_F + \omega_E \cdot U_E + \gamma_{FE} \cdot U_F \cdot U_E + \epsilon_i
\end{align}

The difference identifies:
\begin{equation}
\Delta U_i = \omega_E \cdot U_E + \gamma_{FE} \cdot U_F \cdot U_E
\end{equation}

\textbf{Step 3: Estimation}

With variation in $U_F$ across individuals:
\begin{equation}
\Delta U_i = \alpha + \beta \cdot U_{F,i} + \nu_i
\end{equation}
where $\alpha = \omega_E \cdot \bar{U}_E$ and $\beta = \gamma_{FE} \cdot \bar{U}_E$.

Thus: $\hat{\gamma}_{FE} = \hat{\beta} / \bar{U}_E$.

\paragraph{Result.}

Suppose $\hat{\beta} = 0.15$ (SE = 0.03) and $\bar{U}_E = 0.5$. Then:
\begin{equation}
\hat{\gamma}_{FE} = 0.15 / 0.5 = 0.30 \quad (95\% \text{ CI: } [0.18, 0.42])
\end{equation}

\paragraph{Discussion.}

The positive $\gamma_{FE} = 0.30$ indicates Financial and Emotional utility
are complements: financial gains are more valuable when accompanied by emotional
support. This is consistent with EBF predictions and distinguishes it from
additive models ($\gamma = 0$).


\subsection{Example 2: Learning Dynamics After a Policy Shock}

\paragraph{Setup.}
A government introduces a new carbon tax at $t=0$. We observe how households'
reference structure $C^*$ adjusts over time using panel data with $N = 1,000$
households observed monthly for 24 months.

\paragraph{Model.}
The reference structure follows the learning dynamics from AZ-L1:
\begin{equation}
C_t = C_{t-1} + \alpha \cdot (\Phi(C_{t-1}, \Psi_{new}) - C_{t-1}) + \epsilon_t
\end{equation}

where $\Psi_{new}$ includes the carbon tax.

\paragraph{Application.}

\textbf{Step 1: Estimate adjustment speed $\alpha$}

Regress monthly changes on the gap:
\begin{equation}
\Delta C_t = \alpha \cdot (C^*_{new} - C_{t-1}) + \nu_t
\end{equation}

Using household-level data: $\hat{\alpha} = 0.08$ (SE = 0.01).

\textbf{Step 2: Calculate half-life}
\begin{equation}
t_{1/2} = \frac{\ln(2)}{\alpha} = \frac{0.693}{0.08} \approx 8.7 \text{ months}
\end{equation}

\textbf{Step 3: Compute convergence path}
\begin{equation}
C_t = C^*_{new} + (C_0 - C^*_{new}) \cdot e^{-0.08t}
\end{equation}

\paragraph{Result.}

After 24 months, the gap has closed by:
\begin{equation}
1 - e^{-0.08 \times 24} = 1 - e^{-1.92} \approx 85\%
\end{equation}

\paragraph{Discussion.}

The $\alpha = 0.08$ estimate implies slow adjustment---households take nearly
9 months to close half the gap to the new reference. This has policy implications:
short-term evaluations (e.g., 6 months) will underestimate long-run behavioral
change by approximately 40\%.


\subsection{Example 3: Aggregation with Heterogeneous Thresholds}

\paragraph{Setup.}
A firm wants to predict adoption rates for a sustainable product.
Individual adoption requires $WAX_i \geq \theta_i$, where thresholds
$\theta_i$ are heterogeneous across the population.

\paragraph{Model.}
From AZ-H1, thresholds follow a distribution:
\begin{equation}
\theta_i \sim \text{LogNormal}(\mu_\theta, \sigma_\theta^2)
\end{equation}

with $\mu_\theta = 0.5$ and $\sigma_\theta = 0.3$ (estimated from pilot).

\paragraph{Application.}

\textbf{Step 1: Compute WAX for the product}

Using EBF:
\begin{equation}
WAX = 0.65 \quad \text{(from CORE-READY calculation)}
\end{equation}

\textbf{Step 2: Calculate adoption probability}
\begin{equation}
P(\text{adopt}) = P(\theta_i \leq WAX) = P(\theta_i \leq 0.65)
\end{equation}

For LogNormal:
\begin{equation}
P(\text{adopt}) = \Phi\left(\frac{\ln(0.65) - 0.5}{0.3}\right) = \Phi(-3.10) \approx 0.001
\end{equation}

Wait---this seems too low. Re-check parameters.

\textbf{Step 3: Sensitivity analysis}

If $\mu_\theta = -0.5$ (lower thresholds):
\begin{equation}
P(\text{adopt}) = \Phi\left(\frac{\ln(0.65) - (-0.5)}{0.3}\right) = \Phi(0.23) \approx 0.59
\end{equation}

\paragraph{Result.}

Adoption is highly sensitive to threshold distribution. The mean threshold
$\mathbb{E}[\theta] = e^{\mu + \sigma^2/2}$ determines whether the product
achieves mass adoption or niche uptake.

\paragraph{Discussion.}

This example shows why heterogeneity matters: a single ``representative
threshold'' $\bar{\theta}$ would give binary predictions (adopt/not adopt),
while the distributional approach gives nuanced adoption rates. Segmentation
by $\theta$-type enables targeted interventions.


%%%%%%%%%%%%%%%%%%%%%%%%%%%%%%%%%%%%%%%%%%%%%%%%%%%%%%%%%%%%%%%%%%%%%%%%%%%%%%%
\section{Implications for Practice}
\label{sec:az-practice}
%%%%%%%%%%%%%%%%%%%%%%%%%%%%%%%%%%%%%%%%%%%%%%%%%%%%%%%%%%%%%%%%%%%%%%%%%%%%%%%

\begin{tcolorbox}[colback=green!5!white, colframe=green!75!black,
    title=Practical Implications]
\begin{enumerate}
\item \textbf{For Empirical Researchers:} Use within-person and cross-level
variation to identify EBF parameters. Standard cross-sectional regression
is insufficient due to confounding.

\item \textbf{For Theorists:} Verify axiom consistency before proposing extensions.
The microfoundation chain (P1--P4 $\to$ axioms $\to$ predictions) must remain intact.

\item \textbf{For Policy Makers:} Account for learning dynamics. Policy effects
depend on how quickly agents update reference structures. Immediate evaluation
may underestimate long-run effects.

\item \textbf{For Practitioners:} Recognize heterogeneity. Population-level
interventions work differently for different parameter types. Segment by
$\varphi$ (thinking style) and $\theta$ (action threshold).
\end{enumerate}
\end{tcolorbox}


% =============================================================================
%                    PART 4: BACK MATTER
% =============================================================================

%%%%%%%%%%%%%%%%%%%%%%%%%%%%%%%%%%%%%%%%%%%%%%%%%%%%%%%%%%%%%%%%%%%%%%%%%%%%%%%
\section{Summary}
\label{sec:az-summary}
%%%%%%%%%%%%%%%%%%%%%%%%%%%%%%%%%%%%%%%%%%%%%%%%%%%%%%%%%%%%%%%%%%%%%%%%%%%%%%%

\begin{tcolorbox}[colback=green!10!white, colframe=green!75!black,
    title=\textbf{Appendix AZ: Summary}]

\textbf{Core Question:} What methodological foundations ensure EBF's scientific rigor?

\textbf{Structure:} 9 Parts covering the complete methodological toolkit

\textbf{Main Contributions by Part:}
\begin{itemize}[nosep]
    \item \textbf{Part 1}: Microfoundation from primitives P1--P4 to axiom system
    \item \textbf{Part 2}: Bayesian learning dynamics for $C^*$ convergence
    \item \textbf{Part 3}: Gradient adjustment with stability analysis
    \item \textbf{Part 4}: Distributional heterogeneity and aggregation rules
    \item \textbf{Part 5}: Identification strategies distinguishing EBF
    \item \textbf{Part 6}: Computational methods and algorithms for $C^*$
    \item \textbf{Part 7}: Validation diagnostics and robustness checks
    \item \textbf{Part 8}: Software implementation (pseudocode)
    \item \textbf{Part 9}: Bayesian prior specification and elicitation
\end{itemize}

\textbf{Key Outputs:}
\begin{itemize}[nosep]
    \item \textbf{Axioms}: AZ-M1/M2 (microfoundation), AZ-L1/L2 (learning),
          AZ-D1 (dynamics), AZ-H1 (heterogeneity), AZ-I1 (identification)
    \item \textbf{Algorithms}: AZ-C1/C2/C3 (computational), AZ-S1--S5 (software)
    \item \textbf{Diagnostics}: AZ-V1--V4 (validation protocols)
    \item \textbf{Priors}: Recommended distributions for all EBF parameters
\end{itemize}

\textbf{Worked Examples:} 3 (Complementarity identification, Learning dynamics,
Heterogeneous aggregation)

\textbf{Critical Objections Addressed:} 5 (Circularity, Learning assumptions,
Multiple equilibria, Lucas Critique, External validity)

\textbf{Integration:} Full METHOD ecosystem (AN, R, E, AL), FORMAL-PROOF (D),
CORE-WHERE (BBB), CORE-AWARE/READY (AU, AV), Behavioral Change Journey (Ch. 13--16)

\end{tcolorbox}


%%%%%%%%%%%%%%%%%%%%%%%%%%%%%%%%%%%%%%%%%%%%%%%%%%%%%%%%%%%%%%%%%%%%%%%%%%%%%%%
\section{Glossary of Symbols}
\label{sec:az-glossary}
%%%%%%%%%%%%%%%%%%%%%%%%%%%%%%%%%%%%%%%%%%%%%%%%%%%%%%%%%%%%%%%%%%%%%%%%%%%%%%%

\begin{tcolorbox}[colback=blue!5!white, colframe=blue!75!black]
\textbf{Central Reference:} For complete symbol definitions, see
\textbf{Appendix G} (\textit{Glossary and Notation}, \S\ref{app:glossary}).

This section lists only symbols introduced or specialized in this appendix.
\end{tcolorbox}

\vspace{0.5em}

\begin{table}[htbp]
\centering
\caption{Symbols Introduced in Appendix AZ}
\label{tab:az-symbols}
\renewcommand{\arraystretch}{1.3}
\begin{tabular}{@{}clp{6cm}c@{}}
\toprule
\textbf{Symbol} & \textbf{Name} & \textbf{Definition} & \textbf{Also in G?} \\
\midrule
$\Phi$ & Update operator & $\Phi: \mathcal{C} \times \mathcal{P} \to \mathcal{C}$, maps current state to next & NEW \\
$J$ & Jacobian matrix & $\partial \Phi / \partial C$ at fixed point & NEW \\
$\mathcal{H}_t$ & Information history & All information available to agent at time $t$ & NEW \\
$\alpha_t$ & Learning rate & Stepsize in adaptive learning, $\alpha_t \to 0$ & NEW \\
$\eta$ & Adjustment speed & Speed of gradient adjustment in dynamics & NEW \\
$V(C, \Psi)$ & Potential function & Lyapunov function for stability analysis & NEW \\
$k$ & Cognitive level & Level in hierarchical expectation formation & NEW \\
$F(\Theta | \Psi)$ & Parameter distribution & Population distribution of individual parameters & NEW \\
$\lambda_{\min}$ & Eigenvalue & Smallest eigenvalue of $I - J$ (convergence rate) & NEW \\
\bottomrule
\end{tabular}
\end{table}


%%%%%%%%%%%%%%%%%%%%%%%%%%%%%%%%%%%%%%%%%%%%%%%%%%%%%%%%%%%%%%%%%%%%%%%%%%%%%%%
\section{Critical Foundations and Objections}
\label{sec:az-foundations}
%%%%%%%%%%%%%%%%%%%%%%%%%%%%%%%%%%%%%%%%%%%%%%%%%%%%%%%%%%%%%%%%%%%%%%%%%%%%%%%

\subsection{Objection 1: Microfoundation Circularity}

\textbf{Objection:} The primitives P1--P4 already assume preferences exist.
This is not a true microfoundation but a restatement.

\textbf{Response:} P1--P4 are minimal assumptions that any preference-based theory
must make. The contribution is showing that EBF's complex structure (levels,
dimensions, complementarities, context) follows from these primitives rather than
being imposed ad hoc. The derivation is constructive, not circular.

\subsection{Objection 2: Learning Assumptions Too Strong}

\textbf{Objection:} Bayesian learning requires common knowledge of the model
structure and computational capabilities beyond human capacity.

\textbf{Response:} AZ-L2 explicitly allows for bounded rationality via cognitive
levels $k$. Level-0 (naive) and level-1 (adaptive) learning require no model
knowledge. The Bayesian ideal (level-2) is a limiting case. EBF is compatible
with behavioral learning theories \citep{PAP-camerer2004cognitive}.

\subsection{Objection 3: Multiple Equilibria Undermine Predictions}

\textbf{Objection:} If multiple $C^*$ exist, the model makes no unique predictions.
This is unfalsifiable.

\textbf{Response:} Multiple equilibria are an empirical feature, not a theoretical
weakness. EBF predicts: (1) which equilibria are stable, (2) the basins of
attraction, (3) conditions for regime shifts. These are testable predictions.
The critique applies equally to any model with multiple equilibria (game theory,
macroeconomics).

\subsection{Objection 4: Lucas Critique Applies}

\textbf{Objection:} EBF parameters ($\gamma$, $\omega$, $\lambda$) are not
``deep'' structural parameters. Policy interventions will change them, making
the model useless for counterfactual analysis \citep{lucas1976econometric}.

\textbf{Response:} The Lucas Critique is addressed in three ways:

\begin{enumerate}
\item \textbf{Endogenous $C^*$:} The reference structure $C^* = \Phi(C^*, \Psi)$
explicitly depends on policy environment $\Psi$. When policy changes, $C^*$
adjusts---this is a feature, not a bug.

\item \textbf{Primitive Parameters:} P1--P4 are truly primitive (preference
existence, completeness, transitivity, context-sensitivity). These do not
change with policy.

\item \textbf{Learning Dynamics:} Part 2 specifies how agents update beliefs
about parameters. This makes parameter evolution predictable.
\end{enumerate}

EBF is \textit{Lucas-robust}: it models how behavior changes with policy,
rather than assuming fixed behavioral rules.

\subsection{Objection 5: External Validity Concerns}

\textbf{Objection:} Parameters estimated in one context (e.g., laboratory,
Western sample) may not generalize to other contexts (field, non-Western).

\textbf{Response:} EBF explicitly addresses this via the context framework:

\begin{enumerate}
\item \textbf{Context-Indexed Parameters:} All parameters are functions of $\Psi$:
$\gamma(\Psi)$, $\omega(\Psi)$, etc. Generalization requires specifying how
$\Psi$ differs across contexts.

\item \textbf{Hierarchical Estimation:} Use multilevel models:
\begin{equation}
\theta_{ij} = \bar{\theta} + u_j + v_i \quad \text{where } j \text{ indexes context}
\end{equation}

\item \textbf{Domain Knowledge:} AZ-I1 (Falsifiability) requires observable
implications. External validity is an empirical question, not a theoretical
limitation.
\end{enumerate}

The honest position: EBF parameters \textit{should} vary across contexts.
The framework provides tools to model this variation systematically.


%%%%%%%%%%%%%%%%%%%%%%%%%%%%%%%%%%%%%%%%%%%%%%%%%%%%%%%%%%%%%%%%%%%%%%%%%%%%%%%
\section{Open Issues and Research Agenda}
\label{sec:az-open}
%%%%%%%%%%%%%%%%%%%%%%%%%%%%%%%%%%%%%%%%%%%%%%%%%%%%%%%%%%%%%%%%%%%%%%%%%%%%%%%

\subsection{Methodological Priorities}

\begin{enumerate}
    \item \textbf{Computational Tractability:} Develop efficient algorithms for
    computing $C^*$ in high-dimensional $\Psi$ spaces. Current methods scale poorly
    beyond 8 context dimensions.

    \item \textbf{Finite-Sample Identification:} Establish minimum sample size
    requirements for identifying $\gamma$ parameters with precision. Pilot studies
    suggest $N > 200$ per cell for cross-dimensional complementarity.
\end{enumerate}

\subsection{Empirical Priorities}

\begin{enumerate}
    \item \textbf{Learning Speed Estimation:} Estimate $\eta$ (adjustment speed)
    across domains. Preliminary evidence suggests $\eta_{financial} > \eta_{social}$
    but systematic measurement is lacking.

    \item \textbf{Heterogeneity Structure:} Map the empirical distribution $F(\Theta | \Psi)$
    for key populations. Current work relies on assumed distributions.
\end{enumerate}

\begin{table}[htbp]
\centering
\caption{Research Roadmap for Appendix AZ}
\label{tab:az-roadmap}
\renewcommand{\arraystretch}{1.3}
\begin{tabular}{@{}clcl@{}}
\toprule
\textbf{\#} & \textbf{Issue} & \textbf{Priority} & \textbf{Status} \\
\midrule
1 & Computational efficiency for $C^*$ & HIGH & Active (Part 6 algorithms) \\
2 & Finite-sample identification bounds & HIGH & Preliminary results \\
3 & Learning speed $\eta$ estimation & MEDIUM & Pilot data (Tab.~\ref{tab:az-empirical}) \\
4 & $F(\Theta|\Psi)$ mapping & MEDIUM & Survey design \\
5 & Segment-specific parameters & HIGH & Started (Tab.~\ref{tab:az-segments}) \\
6 & Stage transition probabilities & MEDIUM & Framework ready \\
7 & Cross-validation of empirical estimates & HIGH & Needed \\
8 & Higher-order expectations ($k \geq 3$) & LOW & Theoretical \\
\bottomrule
\end{tabular}
\end{table}


%%%%%%%%%%%%%%%%%%%%%%%%%%%%%%%%%%%%%%%%%%%%%%%%%%%%%%%%%%%%%%%%%%%%%%%%%%%%%%%
\section{References}
\label{sec:az-references}
%%%%%%%%%%%%%%%%%%%%%%%%%%%%%%%%%%%%%%%%%%%%%%%%%%%%%%%%%%%%%%%%%%%%%%%%%%%%%%%

\begin{tcolorbox}[colback=blue!5!white, colframe=blue!75!black]
\textbf{Master Bibliography:} For complete reference details, see
\texttt{bcm2\_master\_references.bib} in the repository.

This section lists appendix-specific references and data sources.
\end{tcolorbox}

\vspace{0.5em}

\subsection{Key References for This Appendix}

The following works are particularly relevant for the content of this appendix:

\begin{itemize}[nosep]
    \item \textbf{Microfoundations:} \citet{lucas1976econometric}, \citet{kirman1992whom}
    \item \textbf{Learning:} \citet{sargent1993bounded}, \citet{evans2001learning}
    \item \textbf{Dynamics:} \citet{benveniste1990adaptive}, \citet{arthur2021foundations}
    \item \textbf{Identification:} \citet{manski1995identification}
    \item \textbf{Bounded rationality:} \citet{stahl1995players}, \citet{nagel1995unraveling},
    \citet{PAP-camerer2004cognitive}
    \item \textbf{Aggregation:} \citet{gorman1953community}
\end{itemize}

\subsection{Appendix-Specific References}

\begin{thebibliography}{99}

\bibitem{arthur2021foundations} Arthur, W.B. (2021). Foundations of complexity economics.
\textit{Nature Reviews Physics}, 3, 136--145.

\bibitem{benveniste1990adaptive} Benveniste, A., Métivier, M., \& Priouret, P. (1990).
\textit{Adaptive Algorithms and Stochastic Approximations}. Springer.

\bibitem{PAP-camerer2004cognitive} Camerer, C.F., Ho, T.H., \& Chong, J.K. (2004).
A cognitive hierarchy model of games. \textit{Quarterly Journal of Economics}, 119(3), 861--898.

\bibitem{evans2001learning} Evans, G.W., \& Honkapohja, S. (2001).
\textit{Learning and Expectations in Macroeconomics}. Princeton University Press.

\bibitem{gorman1953community} Gorman, W.M. (1953). Community preference fields.
\textit{Econometrica}, 21(1), 63--80.

\bibitem{kirman1992whom} Kirman, A.P. (1992). Whom or what does the representative individual represent?
\textit{Journal of Economic Perspectives}, 6(2), 117--136.

\bibitem{lucas1976econometric} Lucas, R.E. (1976). Econometric policy evaluation: A critique.
\textit{Carnegie-Rochester Conference Series on Public Policy}, 1, 19--46.

\bibitem{manski1995identification} Manski, C.F. (1995). \textit{Identification Problems in the Social Sciences}.
Harvard University Press.

\bibitem{nagel1995unraveling} Nagel, R. (1995). Unraveling in guessing games: An experimental study.
\textit{American Economic Review}, 85(5), 1313--1326.

\bibitem{sargent1993bounded} Sargent, T.J. (1993). \textit{Bounded Rationality in Macroeconomics}.
Oxford University Press.

\bibitem{stahl1995players} Stahl, D.O., \& Wilson, P.W. (1995). On players' models of other players:
Theory and experimental evidence. \textit{Games and Economic Behavior}, 10(1), 218--254.

\end{thebibliography}

\subsection{Data Sources}

\begin{itemize}[nosep]
    \item \textbf{Parameter Estimates (Table~\ref{tab:az-empirical}):} Synthesized from
    published meta-analyses; see individual citations in table.
    \item \textbf{Convergence Benchmarks (Table~\ref{tab:az-complexity}):} Computational
    experiments on standard hardware (Intel i7, 16GB RAM).
    \item \textbf{Segment Parameters (Table~\ref{tab:az-segments}):} Illustrative values
    based on behavioral segmentation literature; require empirical validation.
\end{itemize}

\noindent\textit{Note:} No primary data collection was performed for this methodological
appendix. All empirical estimates are secondary, drawn from cited literature.


%%%%%%%%%%%%%%%%%%%%%%%%%%%%%%%%%%%%%%%%%%%%%%%%%%%%%%%%%%%%%%%%%%%%%%%%%%%%%%%
% CROSS-REFERENCE FOOTER (Required)
%%%%%%%%%%%%%%%%%%%%%%%%%%%%%%%%%%%%%%%%%%%%%%%%%%%%%%%%%%%%%%%%%%%%%%%%%%%%%%%

\vspace{1em}
\hrule
\vspace{0.5em}

\noindent\textit{Cross-references:}
Appendix G (Central Glossary),
Appendix A (FORMAL-DERIVE: Mathematical Foundations),
Appendix D (FORMAL-PROOF: Theorem Demonstrations),
Appendix T (REF-META: Metatheory),
Appendix BBB (CORE-WHERE: Parameter Estimation),
Appendix AN (METHOD-LLMMC: LLM Monte Carlo).

\noindent\textit{Master Bibliography:} \texttt{bcm2\_master\_references.bib}

% =============================================================================
% END OF APPENDIX AZ
% =============================================================================
